\section*{Limitations}
While our AgentHallu marks a critical advancement in automated hallucination attribution for LLM-based agents, it is important to recognize several limitations. First, although AgentHallu spans 5 primary categories and 14 subcategories, it remains challenging to fully anticipate and represent emerging hallucination patterns as agent frameworks, tool ecosystems, and interaction protocols rapidly evolve. Therefore, the dataset should be continuously expanded to keep pace with new agent capabilities. Second, AgentHallu primarily targets text-based trajectories and does not consider multimodal agent settings grounded in images, audio, or other modalities. Given the growing adoption of multimodal agents, future work should explore extending the attribution framework to encompass these broader multimodal interactions.

% Third, the current benchmark is predominantly English, which limits its applicability to multilingual agent deployments; future work should expand coverage to additional languages and evaluate whether attribution behaviors and error types transfer across linguistic contexts.

\section*{Ethical Considerations}
AgentHallu is strictly free of personally identifiable information and offensive content. The benchmark is exclusively sourced from publicly accessible datasets and repositories, as well agent trajectories generated under controlled settings, explicitly avoiding sensitive or restricted data sources. Designed for academic research, AgentHallu focuses on enhancing the reliability of autonomous agents. Through adherence to strict data integrity protocols and ethical standards, AgentHallu establishes a responsible foundation for the automated attribution of agent hallucinations.